\newcommand{\ModuleID}{EL-III.8}
\newcommand{\ModuleTitle}{Prefaces and the Roman Canon}
\newcommand{\ModuleStage}{Stage III — Independent Missal Reading}
\newcommand{\ModuleUnit}{III.8}
\newcommand{\ModuleOutcome}{Hold a long liturgical period together through its finite spine, infinitive frame, relative chains, object series, and destinations}
\newcommand{\ModulePrerequisites}{EL-III.7 and accurate reconstruction of compressed order}
\newcommand{\ModuleExtensionPractice}{Worksheets III.30 and III.31 remain optional Christmas-Preface and Canon-opening practice}
\newcommand{\ModuleNext}{EL-III.9, The Temporal Cycle}
\newcommand{\ModuleMastery}{On an unfamiliar Preface or Canon period, identify the main completion, attach all subordinate and non-finite material, resolve reference and agreement chains, and explain the function of one delayed or repeated member.}
\newcommand{\ModuleDelayNote}{}
\newcommand{\ModuleLessonSource}{\CourseRoot/shared/content/volume3/all}
\newcommand{\ModuleExerciseSource}{\CourseRoot/shared/exercises/volume3_4/volume3}
\newcommand{\ModuleWorkedBank}{III}
\newcommand{\ModuleExerciseBank}{III}
\newcommand{\ModuleWorksheetVariants}{A,D}
\newcommand{\ModuleScopeNote}{The packet is linguistic analysis of selected fixed texts, not ceremonial instruction or a theological exposition of the Canon.}
\newcommand{\ModulePractice}{%
  \SubmoduleExtension{The Preface frame}
    {The opening predicates anticipate an act that completes their sense.}
    {Copy the lesson's Preface frame, box the predicate series, bracket the completing infinitive unit, and show where its subject is expressed or understood.}
    {Why are the opening neuter predicates not evidence for a masculine subject?}
    {The model keeps the neuter predicate adjectives with the anticipated act or infinitive construction and identifies the copula. No masculine noun may be invented from English idiom.}
  \SubmoduleExtension{The opening petition of the Canon}
    {Repeated objects and relative links form one governed petition despite intervening modifiers.}
    {Draw the petition's finite spine, then attach every object, apposition, adjective, relative, and recipient in dependency order.}
    {Name the single governing petition before listing its members.}
    {A complete diagram begins with the governing verbal center and shows the repeated accusatives as its coordinated or expanding object series; each relative has an explicit antecedent.}
  \SubmoduleExtension{Predicate adjectives in the \emph{Quam oblationem}}
    {A string of agreeing accusatives can describe the object predicatively rather than add new objects.}
    {Use the lesson passage to separate object phrase, governing infinitive, and predicate accusatives; then change the project-composed object to plural and cascade all agreement.}
    {State the test distinguishing a predicate accusative from an independent object.}
    {The adjectives agree with and are predicated of the same object through the governing construction. In the transformation every linked adjective changes with the object; no new verbal object is introduced.}
  \SubmoduleExtension{Command plus passive infinitive}
    {A command may govern an infinitive whose patient and destination must be traced separately.}
    {On the lesson example, label command, infinitive, logical subject or patient, personal agent if any, and both destination phrases.}
    {Explain why two occurrences of \Latin{in} need not govern the same case.}
    {The model attaches the passive infinitive to the command and assigns each \Latin{in}-phrase from motion or location and its printed case. A single English preposition is not a sufficient parse.}
  \SubmoduleExtension{Canon families}
    {A family inventory reduces cognitive load without turning formula into one undifferentiated gloss.}
    {Create four construction-aware families from the lesson: petition, offering, acceptance, and commemoration; give one lemma and one governed form for each.}
    {Retrieve the construction, not merely the English label, for two entries.}
    {Each accepted entry cites an actual lesson form and names its object, complement, or case government. Entries with only English thematic labels do not pass.}
}
